\section{Limitations: de Benchmark Capability a Software Reliability}

La clase no terminó en la main figure, sino que mostró la brittle completion, la relación entre model-size/compute y las conclusiones. Un pass@k alto puede demostrar que existen muestras funcionalmente correctas en un benchmark estrecho, pero no cubre automáticamente specification, security, composition ni mantenimiento.

\subsection{Brittleness y composition}

Esta sección revisa primero los failures y después explica los límites de la main figure. La prompt brittleness y la degradación en composition muestran que un benchmark de funciones cortas no puede representar el software a repository-scale.

\lecturefigure{slide-44-limitations.jpg}{Las limitations de Codex incluyen la prompt brittleness y la degradación rápida en composition más compleja.}{Video oficial, 46:53.}

\begin{knowledgebox}{Lectura de la figura: a la izquierda, fragilidad local; a la derecha, tendencia con la escala}
Un cambio pequeño en el prompt puede provocar una completion errónea; a medida que la tarea requiere combinar más pasos, la exactitud disminuye. El tamaño del modelo mejora la curva, pero no garantiza la extrapolación a programas largos ni a repositories reales. Conviene evaluar paraphrase, hidden requirements, multi-file dependency, security y regression.
\end{knowledgebox}

\begin{warningbox}{La functional correctness todavía no es production readiness}
El código de producción también necesita readability, types, security, performance, licenses, dependency policy, observabilidad y human review. Aunque todas las unit tests pasen, pueden existir entradas no cubiertas o efectos secundarios peligrosos. La salida del modelo debe integrarse al proceso de ingeniería de software, no saltárselo.
\end{warningbox}

\subsection{Los cuatro juicios duraderos de la slide de conclusiones}

Al final, la slide de conclusiones se reformula como proposiciones verificables: bajo qué condiciones de data, metric y cost se cumplen el modelado generativo, la scale, la transferencia entre modalidades y el sampling. Al repasar toda la sesión, conviene colocar sample quality, representation transfer, benchmark pass y deployment reliability en columnas distintas, para no dejar que una sola historia de éxito cubra todos los niveles.

\lecturefigure{slide-45-conclusion.jpg}{El ponente resume toda la sesión con el preentrenamiento generativo, la scale, la cross-modality y el sampling.}{Video oficial, 47:38.}

\begin{knowledgebox}{Lectura de la figura: reformular las conclusiones como proposiciones verificables}
El modelado generativo puede aprovechar datos sin etiquetar; la scale suele mejorar el sample y la transfer; una misma recipe puede transferirse a image/code; el sampling puede amplificar los rare success de la distribución. Cada afirmación debe acotarse en data, model, metric y cost. La extensión más importante es esta: cuando existe un verifier, el modelo generativo puede formar un sistema junto con la búsqueda; cuando el verifier no es confiable, el muestreo múltiple también puede amplificar el riesgo.
\end{knowledgebox}

\teachervoice{Al cierre, el ponente enfatiza a la vez la amplitud del generative modeling y sus limitations. Recordatorio de la clase: lo destacable de Codex no es que «los problemas de código estén resueltos», sino que el code ofrece un entorno de investigación que se puede ejecutar, muestrear y verificar, en el que la distribución del modelo y la búsqueda del sistema pueden medirse por separado.}

\subsection{Resumen de la sección}

La benchmark capability es evidencia necesaria, no confiabilidad del software. Codex nos permite ver la estructura del sistema de generación, muestreo, prueba y selección, y también deja más claros la prompt brittleness, la composition y los límites de seguridad.
